% capitulo4.tex
\chapter{Marco Teórico y Referencial}

El presente capítulo establece el sustento teórico y referencial del trabajo. Se parte de los
antecedentes histórico-normativos que condicionaron la evolución del diseño de placas base en
México, se describen los mecanismos de transferencia de fuerzas y los modos de falla relevantes,
y se documentan con precisión las deficiencias de las NTC-DCEA 2017 y los cambios sustanciales
introducidos por las NTC-DCEA 2023. Las secciones de análisis comparativo cuantitativo y
discusión de resultados se incorporarán en las entregas subsecuentes del trabajo.

% ============================================================
% SECCIÓN 4.1 — CONTEXTO GEOTECTÓNICO Y MARCO NORMATIVO
% ============================================================
\section{Contexto Geotectónico y Marco Normativo en México}

La República Mexicana ocupa una posición geotectónica de singular complejidad: su territorio es
escenario de la convergencia de la placa de Norteamérica con las placas de Cocos y Rivera en la
franja costera del Pacífico, así como de la interacción con la placa del Caribe en la región
sureste del país. Esta confluencia de placas genera una actividad sísmica continua e intensa,
con registros históricos que incluyen algunos de los eventos más destructivos del siglo XX en el
continente americano \cite{UNAM_SSN_2017}.

La Ciudad de México suma a esta condición sísmica de base un factor de amplificación dinámica
excepcional: gran parte de su área metropolitana está asentada sobre los sedimentos lacustres
del antiguo lago de Texcoco, compuestos por arcillas de alta plasticidad con velocidades de onda
cortante de apenas 40 a 80~m/s en sus estratos más blandos. Este perfil estratigráfico genera
amplificaciones del movimiento del suelo de hasta un orden de magnitud respecto a la roca firme,
con periodos predominantes del sitio que pueden alcanzar de 2 a 4~s en las zonas lacustres más
profundas \cite{UNAM_SSN_2017}.

La evolución del marco normativo sismorresistente en México ha estado históricamente condicionada
por los grandes sismos:

\begin{itemize}
  \item El sismo de 1985 ($M_s~8.1$, epicentro en Michoacán) causó la destrucción de más de 800
    edificaciones en la Ciudad de México y aceleró la publicación de las Normas Técnicas
    Complementarias de 1987, que introdujeron criterios de microzonificación sísmica, espectros de
    diseño diferenciados por tipo de suelo y requerimientos de compacidad para elementos de acero.

  \item Las revisiones de 1993 y 2004 incorporaron refinamientos en el cálculo de derivas de
    entrepiso y en los estados límite de servicio, alineando parcialmente la práctica mexicana con
    los desarrollos del American Institute of Steel Construction \cite{AISC_360_2022}.

  \item Las NTC-DCEA 2017 representaron un avance conceptual al incluir en su Apéndice~B las
    primeras disposiciones específicas para el diseño de placas base. No obstante, su contenido
    técnico fue redactado y cerrado por los comités técnicos con anterioridad al sismo del
    19~de septiembre de 2017, por lo que no pudo incorporar las lecciones mecánicas de dicho
    evento \cite{GCDMX_NTC_DCEA_2017}.

  \item Las NTC-DCEA 2023 constituyen la primera edición que incorpora explícitamente las lecciones
    del 19-S~2017, elevando el diseño de placas base a la categoría de capítulo principal mandatorio
    (Capítulo~13) e integrando los principios del Diseño por Capacidad para las conexiones
    columna-cimentación \cite{GCDMX_NTC_DCEA_2023}.
\end{itemize}

% ============================================================
% SECCIÓN 4.2 — EL SISMO DEL 19-S 2017 Y LA EVIDENCIA FORENSE
% ============================================================
\section{El Sismo del 19 de Septiembre de 2017 y la Evidencia Forense Estructural}

El 19 de septiembre de 2017 a las 13:14~h (hora local), un sismo de magnitud momento $M_w~7.1$
sacudió la región centro-sur de México, con epicentro en los límites entre los estados de Puebla
y Morelos, a escasos 120~km de la Ciudad de México \cite{SSN_Puebla_2017}. A diferencia del
sismo de 1985, este evento fue de origen intraplaca (no de subducción), lo que le confirió un
contenido frecuencial caracterizado por periodos predominantes cortos e intermedios (0.1--1.5~s),
particularmente agresivos para estructuras de 4 a 10 niveles. El balance de daños fue grave:
múltiples colapsos parciales o totales en edificios de uso habitacional y comercial, con víctimas
fatales y pérdidas materiales cuantiosas.

La \textsc{SMIE} organizó brigadas de inspección forense que documentaron patrones de falla
recurrentes en conexiones columna-cimentación de estructuras de acero \cite{SMIE_2017}:

\begin{itemize}
  \item \textbf{Arrancamiento de cono de concreto}: identificado como el modo de falla más
    frecuente y más frágil. Los pernos de anclaje extrajeron el concreto del pedestal sin que
    la columna metálica hubiera experimentado deformaciones plásticas significativas, lo que
    evidenció una violación del principio de jerarquía de resistencias del Diseño por Capacidad.
  \item \textbf{Fluencia prematura de la placa base}: en casos donde los anclas sobrevivieron, la
    placa base sufrió deformaciones permanentes que comprometieron la transmisión uniforme de
    cargas verticales al pedestal, generando aplastamientos locales del grout y del concreto.
  \item \textbf{Aplastamiento del mortero de nivelación}: consecuencia de la concentración de
    presiones de contacto bajo el patín de compresión de la columna, especialmente en placas
    subdimensionadas respecto a la demanda real.
\end{itemize}

Las investigaciones de Tapia~Hernández \cite{TapiaHernandez_2018} aportaron una explicación
complementaria a estos fallos: los factores de sobreresistencia reales del acero estructural
mexicano ($\Omega_o$) superaban consistentemente los valores nominales empleados en el diseño
elástico. En consecuencia, las fuerzas que efectivamente llegaban a las conexiones durante el
sismo superaban las previstas por el análisis espectral, desbordando la capacidad de los sistemas
de anclaje diseñados con la metodología de 2017.

% ============================================================
% SECCIÓN 4.3 — MECANISMOS DE TRANSFERENCIA Y MODOS DE FALLA
% ============================================================
\section{Mecanismos de Transferencia de Fuerzas y Modos de Falla Típicos}

La comprensión de los mecanismos de transferencia de fuerzas en el sistema placa base-anclaje-dado
es prerequisito para el análisis normativo comparativo.

\subsection*{Transferencia bajo compresión pura y pequeña excentricidad}

Cuando la excentricidad de la carga resultante $e = M_u/P_u$ es menor que la excentricidad crítica
$e_{crit}$, toda la placa permanece en contacto con el concreto y el mecanismo de transferencia
predominante es la compresión directa. El estado límite relevante es el aplastamiento del concreto,
verificado a través del esfuerzo de contacto máximo $f_{p,\max}$ comparado con la resistencia
de diseño a compresión amplificada por el factor de confinamiento de área adyacente
\cite{ACI318_2019}.

\subsection*{Transferencia bajo flexocompresión con gran excentricidad}

Cuando $e > e_{crit}$, la placa experimenta levantamiento parcial: parte de ella pierde contacto
con el concreto y los pernos de anclaje del lado en tensión son traccionados. El equilibrio del
sistema se establece entre el bloque equivalente de compresión del concreto ---de profundidad $Y$,
modelado mediante la hipótesis de bloque rectangular de Whitney--- y la fuerza de tensión $T_{ua}$
en los pernos de anclaje \cite{FisherKloiber_2006}. Este es el caso más crítico desde la
perspectiva sísmica, ya que la falla frágil de los anclajes en este régimen puede ser catastrófica.

\subsection*{Modos de falla en el sistema de anclaje}

La clasificación de los estados límite de falla adoptada por las NTC-DCEA 2023 en concordancia
con el Método de Diseño de Capacidad de Concreto (CCD, por sus siglas en inglés) del ACI~318
\cite{ACI318_2019} comprende los siguientes modos principales, ilustrados en las
Figuras~\ref{fig:falla_pullout} y~\ref{fig:falla_cono}:

\begin{enumerate}
  \item \textbf{Fluencia y fractura del acero del perno}: modo dúctil preferente; la fuerza de
    tensión $T_{ua}$ excede la resistencia de fluencia nominal del acero del ancla.
  \item \textbf{Arrancamiento de cono de concreto (\textit{breakout})}: modo frágil; el anclaje
    extrae un volumen de concreto en forma de cono truncado cuya área proyectada $A_{Nc}$
    determina la resistencia nominal $N_{cb}$.
  \item \textbf{Extracción directa (\textit{pull-out})}: falla de adherencia entre el perno y el
    concreto sin formación de cono completo; asociada a anclas sin elemento de cabeza o con
    empotramiento insuficiente.
  \item \textbf{Estallido lateral (\textit{side-face blowout})}: falla del recubrimiento lateral
    del concreto en anclas de cabeza con distancia al borde muy reducida.
  \item \textbf{Falla por pryout (\textit{pryout failure})}: deslizamiento lateral del anclaje
    bajo carga cortante, predominante en anclas cortas.
\end{enumerate}

\begin{figure}[H]
  \centering
  \begin{minipage}[b]{0.46\textwidth}
    \centering
    \includegraphics[width=\textwidth]{\detokenize{figuras/006_Planta de placa base con perfil W al centro..png}}
    \caption{Mecanismo de falla por extracción directa (\textit{pull-out}): pérdida de
      adherencia sin formación de cono de concreto \cite{INDEX_nodate}.}
    \label{fig:falla_pullout}
  \end{minipage}
  \hfill
  \begin{minipage}[b]{0.46\textwidth}
    \centering
    \includegraphics[width=\textwidth]{\detokenize{figuras/007_Fallo por cono de concreto..png}}
    \caption{Mecanismo de falla por arrancamiento de cono de concreto (\textit{breakout}):
      modo frágil predominante en sismos severos \cite{INDEX_nodate}.}
    \label{fig:falla_cono}
  \end{minipage}
\end{figure}

% ============================================================
% SECCIÓN 4.4 — DEFICIENCIAS NTC-DCEA 2017
% ============================================================
\section{Deficiencias del Apéndice B de las NTC-DCEA 2017}

El diseño de placas base en las NTC-DCEA 2017 estaba regulado por el Apéndice~B de dicha norma,
cuya ubicación como apéndice generaba ambigüedad interpretativa respecto a su carácter mandatorio
o complementario \cite{GCDMX_NTC_DCEA_2017}. Esta condición periférica se acompañaba de
limitaciones técnicas sustanciales:

\begin{enumerate}
  \item \textbf{Ausencia de los factores de confinamiento $R_{et}$ y $R_{ev}$}: la metodología de
    2017 no reconocía analíticamente la contribución del refuerzo transversal del dado a la
    resistencia del cono de concreto. Esta omisión conducía a diseños sistemáticamente conservadores
    en dados con refuerzo transversal adecuado y, paradójicamente, no incentivaba la provisión de
    dicho refuerzo como estrategia de mejora de la ductilidad de la conexión. La distribución de
    presiones resultante del paradigma de placa rígida de la NTC-2017 se ilustra en la
    Figura~\ref{fig:presiones_lineales}.

  \item \textbf{Distancia mínima al borde de $6\,d_o \geq 100\ \mathrm{mm}$}: el requisito de la
    NTC-2017 era menos restrictivo que el establecido posteriormente para anclas de alta resistencia,
    lo que podía incrementar la probabilidad del modo de falla por estallido lateral en pedestales
    de dimensiones compactas.

  \item \textbf{Ambigüedad en el Diseño por Capacidad}: la norma de 2017 no incluía lenguaje
    normativo suficientemente explícito que obligara al ingeniero a diseñar las conexiones de base
    para el momento plástico máximo esperado de la columna. Esta ambigüedad permitía legalmente
    diseñar con las reacciones espectrales reducidas por el factor de comportamiento sísmico $Q$,
    lo que resultaba en anclas insuficientes bajo demanda sísmica real \cite{AISC_341_2022}. El
    efecto de palanca (\textit{prying action}), que agrava esta insuficiencia, se ilustra en la
    Figura~\ref{fig:prying_action}.

  \item \textbf{Ausencia de criterios para secciones tubulares y cajón}: las disposiciones del
    Apéndice~B estaban calibradas exclusivamente para perfiles tipo~$I$ (patines anchos), sin
    contemplar las particularidades geométricas de las columnas tubulares y de sección cajón,
    ampliamente utilizadas en la práctica constructiva mexicana.
\end{enumerate}

\begin{figure}[H]
  \centering
  \begin{minipage}[b]{0.47\textwidth}
    \centering
    \includegraphics[width=\textwidth]{\detokenize{figuras/003_Diagrama de distribución lineal de presiones en placa base..png}}
    \caption{Distribución de presiones de contacto bajo la hipótesis de placa rígida, paradigma
      del diseño bajo NTC-DCEA 2017 \cite{Hilti_nodate}.}
    \label{fig:presiones_lineales}
  \end{minipage}
  \hfill
  \begin{minipage}[b]{0.47\textwidth}
    \centering
    \includegraphics[width=\textwidth]{\detokenize{figuras/004_Diagrama esquemático del efecto prying action; placa deformada..png}}
    \caption{Efecto de palanca (\textit{prying action}) en placa base flexible: la flexibilidad
      real incrementa la demanda en los pernos por encima de la solución estática directa
      \cite{Hilti_nodate}.}
    \label{fig:prying_action}
  \end{minipage}
\end{figure}

% ============================================================
% SECCIÓN 4.5 — NTC-DCEA 2023 COMO RESPUESTA TÉCNICA
% ============================================================
\section{Las NTC-DCEA 2023 como Respuesta Técnica}

Las NTC-DCEA 2023 representan un salto cualitativo en la regulación de conexiones
columna-cimentación en México \cite{GCDMX_NTC_DCEA_2023}. Sus cambios principales se describen
a continuación:

\begin{enumerate}
  \item \textbf{Jerarquización normativa}: el diseño de placas base fue elevado del Apéndice~B al
    Capítulo~13 principal, con carácter mandatorio. Se crea adicionalmente el Capítulo~14 para
    conexiones de columnas de sección tubular y cajón, subsanando el vacío de la norma anterior.

  \item \textbf{Diseño por Capacidad (\textit{Capacity Design})}: la NTC-2023 establece que la
    placa base, los pernos de anclaje y el dado de concreto constituyen una ``conexión infalible''
    cuya capacidad de resistencia debe ser superior al momento máximo probable, al cortante
    inelástico máximo y a las tracciones derivadas de la formación de una rótula plástica en la
    base de la columna, evaluados bajo los factores de sobreresistencia $\Omega_o$ del acero
    \cite{AISC_341_2022}.

  \item \textbf{Factor de confinamiento en tensión $R_{et}$}: cuantifica la contribución de los
    estribos del dado a la resistencia al cono de concreto bajo demanda de tensión. Su formulación
    reconoce que los estribos interceptan geométricamente la superficie del cono hipotético de
    falla, restringiendo la propagación frágil de la grieta:
    \begin{equation}
      R_{et} = N_{est} \cdot F_R \cdot \left(0.5 \cdot A_e \cdot \sqrt{f'_c \cdot E_c}\right)
    \end{equation}
    donde $N_{est} = \lfloor h_{ef}/s_{est} \rfloor$ es el número de estribos que interceptan el
    plano del cono hipotético, $F_R$ es el factor de resistencia del material, $A_e$ es el área
    efectiva del estribo en la zona de intersección, $f'_c$ es la resistencia a la compresión del
    concreto y $E_c$ es su módulo de elasticidad.

  \item \textbf{Factor de confinamiento en cortante $R_{ev}$}: cuantifica la contribución de los
    estribos a la resistencia al deslizamiento lateral (cono por cortante):
    \begin{equation}
      R_{ev} = N_{est} \cdot F_R \cdot A_e \cdot F_y
    \end{equation}
    donde $F_y$ es el esfuerzo de fluencia del acero del estribo.

  \item \textbf{Distancia mínima al borde}: modificada a $7\,d_o \geq 100\ \mathrm{mm}$ para
    anclas fabricadas con aceros de alta resistencia A325/A449, lo que representa un incremento del
    16.7\,\% respecto al requisito de la norma de 2017 para este tipo de material.

  \item \textbf{Sinergia con NTC-Concreto 2023}: la metodología de cálculo de resistencia al cono
    de concreto queda integrada con el Capítulo~14 de las Normas Técnicas Complementarias para el
    Diseño y Construcción de Estructuras de Concreto 2023, adoptando formalmente el Método CCD
    \cite{GCDMX_2023_concrete}.
\end{enumerate}

La Figura~\ref{fig:despiece_isometrico} ilustra el despiece tridimensional de la conexión en el
contexto del enfoque integral de diseño propuesto por la NTC-2023.

\begin{figure}[H]
  \centering
  \includegraphics[width=0.60\textwidth]{\detokenize{figuras/008_Despiece isométrico de la conexión; columna W, placa base, pernos de anclaje con tuerca y arandela, capa de grout y pedestal de concreto..png}}
  \caption{Despiece isométrico de la conexión columna-placa base: columna IR, placa de acero,
    pernos de anclaje con tuerca y arandela de alta resistencia, capa de grout y pedestal de
    concreto reforzado con estribos \cite{Shi_2025}.}
  \label{fig:despiece_isometrico}
\end{figure}

% ============================================================
% SECCIÓN 4.6 — BRECHA DE INVESTIGACIÓN
% ============================================================
\section{Brecha de Investigación}

La revisión del estado del arte en el contexto de la normativa mexicana vigente evidencia la
existencia de tres brechas de conocimiento específicas que motivan la presente investigación:

\begin{enumerate}
  \item No existe, a la fecha de elaboración de este trabajo, un análisis paramétrico publicado
    que cuantifique el impacto numérico de los factores $R_{et}$ y $R_{ev}$ de la NTC-DCEA 2023
    sobre la resistencia al cono de concreto, en función de variables prácticas de diseño como la
    profundidad de empotramiento $h_{ef}$, el diámetro del ancla $d_o$ y el espaciado de los
    estribos $s_{est}$.

  \item No se ha documentado sistemáticamente si el cambio en el criterio de distancia mínima al
    borde ($6\,d_o \to 7\,d_o$) tiene un impacto material sobre las dimensiones en planta del
    dado de concreto en el rango de diámetros de ancla más utilizados en la práctica estructural
    mexicana.

  \item No existe una evaluación cuantitativa de la evolución del espesor requerido $t_{req}$ de
    la placa base como función del momento último $M_u$ que compare explícitamente ambas normas
    bajo las mismas condiciones de carga, determinando si los cambios normativos de la NTC-2023
    son o no sensibles a este parámetro de diseño.
\end{enumerate}

La presente tesina aborda estas tres brechas mediante la implementación de algoritmos de diseño
parametrizados en \textsc{Python}, la ejecución de cuatro barridos paramétricos y el análisis
comparativo de los resultados obtenidos bajo ambos marcos normativos. El caso semilla utilizado
corresponde al perfil IR~305$\times$74.5 con placa base de 50$\times$50~cm, concreto con
$f'_c = 250\ \mathrm{kg/cm^2}$ y un régimen de flexocompresión con gran excentricidad
($P_u = 10{,}000\ \mathrm{kg}$, $M_u = 1{,}400{,}000\ \mathrm{kg{\cdot}cm}$, correspondiente a
$e = 140\ \mathrm{cm} \gg e_{crit}$), garantizando la activación de los pernos de anclaje en
tensión y el levantamiento parcial de la placa.
